\chapter{Broken Nose (Nasal Fracture)}
\index{Broken Nose (Nasal Fracture)}

\section*{Definition and Overview}
A broken nose, or nasal fracture, is a crack or break in the bone and cartilage that forms the structure of the nose. It is the most common fracture of the face and is usually caused by a direct blow, such as in a fall, a sporting collision, an assault, or a motor vehicle accident. The nose is composed of thin bones in its upper third and cartilage in its lower portion, which makes it relatively fragile and prone to injury. Most nasal fractures are simple and heal without treatment, but some produce a lasting change in the shape of the nose, persistent nasal obstruction, or an associated head or facial injury that requires more urgent management. Because swelling can hide a developing deformity, timely assessment is important.

\section*{Epidemiology}
Nasal fractures are the most common facial fractures and one of the most frequent reasons for emergency department attendance in the head-and-neck region. They occur most often in teenagers and young adults and are roughly twice as common in males as in females. Leading causes include contact sports (football, boxing, hockey, and martial arts), falls, fights and assaults, and motor vehicle accidents. In young children the nasal bones are softer, smaller, and less prominent, so nasal fractures are less common than in older children and adults. Many minor injuries are never formally diagnosed or treated, so the true incidence is higher than official figures suggest.

\section*{Aetiology and Pathophysiology}
\begin{itemize}
    \item \textbf{Direct blow:} a strike to the front or side of the nose from a fist, elbow, ball, fall, or impact in a vehicle
    \item \textbf{Mechanism of injury:} the thin nasal bones snap under force, while the septal cartilage may bend, split, or be avulsed from its bony attachments
    \item \textbf{Lateral force:} drives the nose sideways, causing deviation and typically injuring the nasal septum, the central dividing wall
    \item \textbf{Frontal force:} can push the nasal bones backward, flattening the bridge and causing telescoping injuries
    \item \textbf{Rich blood supply:} the highly vascular nasal tissues bleed readily, so nosebleeds are common after injury
    \item \textbf{Associated injuries:} septal haematoma, black eyes, and fractures of the cheekbone, eye socket, or upper jaw can accompany a broken nose
\end{itemize}

\section*{Clinical Features}
\begin{itemize}
    \item Pain and tenderness over the nose that worsen on light touch
    \item Rapid swelling of the nose, which can obscure an underlying deformity in the first days
    \item Bruising around the nose and beneath the eyes (raccoon-like black eyes)
    \item A nosebleed, which may be heavy or prolonged
    \item An obvious change in shape: the nose may look crooked, flattened, or pushed to one side
    \item A blocked feeling or difficulty breathing through the nose
    \item A crackling sensation (crepitus) when the nose is gently palpated
    \item Double vision, pain around the eyes, or restricted eye movement, which suggest a more serious facial injury
    \item Clear or watery fluid leaking from the nose after significant head trauma, indicating a possible skull-base fracture
\end{itemize}

\section*{Differential Diagnosis}
\begin{itemize}
    \item \textbf{Soft-tissue injury:} bruising and swelling without a fracture, which can appear identical in the first days
    \item \textbf{Septal haematoma:} a blood collection within the septum that blocks the nose and threatens septal damage if not drained
    \item \textbf{Isolated septal fracture or deviation:} injury to the central cartilage causing obstruction even when the bony nose is intact
    \item \textbf{Facial fractures:} fractures of the cheekbone, orbital floor, or upper jaw that may mimic or coexist with a nasal fracture
    \item \textbf{Skull-base fracture:} cerebrospinal fluid leakage or blood from the nose after major head trauma
    \item \textbf{Nasal foreign body:} in children, a retained object can cause swelling, discharge, and nasal blockage
\end{itemize}

\section*{When to Seek Medical Care}
See a doctor promptly after any nasal injury that is painful, swollen, or bleeding, or if you believe the shape of the nose has changed. Early assessment matters because swelling can hide a deformity that is most easily corrected within the first one to two weeks.

\textbf{Contact your local emergency services for:}
\begin{itemize}
    \item A nosebleed that does not stop after 15--20 minutes of firm, continuous pressure
    \item A nose that appears severely bent, crushed, or flattened, or difficulty breathing through the nose
    \item A growing, painful swelling inside the nose that blocks it (possible septal haematoma)
    \item Clear or watery fluid leaking from the nose after a head injury
    \item Double vision, eye pain, or an inability to move the eye normally
    \item A head injury accompanied by confusion, vomiting, or loss of consciousness
\end{itemize}

\section*{Diagnosis and Investigations}
\begin{itemize}
    \item \textbf{History:} the mechanism of injury, whether there was any loss of consciousness, and whether bleeding has stopped
    \item \textbf{Examination:} inspection for swelling and deformity, gentle palpation for tenderness and crepitus, and examination of the nasal passages for a septal haematoma and septal deviation
    \item \textbf{Assessment of surrounding structures:} checking the eyes, teeth, cheekbones, and bite to identify associated fractures
    \item \textbf{Imaging:} plain X-rays and CT scans are not routinely required for a simple nasal fracture; CT is reserved for suspected complex facial fractures or head injury
    \item \textbf{Follow-up review:} reassessment after swelling has settled, usually at about one week, to decide whether manipulation of the bones is needed
\end{itemize}

\section*{Management}
\begin{itemize}
    \item \textbf{Immediate first aid:} sit leaning forward and pinch the soft part of the nose firmly for 15--20 minutes to stop bleeding; apply a cloth-wrapped ice pack to reduce swelling
    \item \textbf{Pain relief:} simple analgesia such as paracetamol or ibuprofen for the first few days
    \item \textbf{Minor fractures:} most simple, undisplaced breaks require no treatment; the nose is left to heal and swelling subsides over one to two weeks
    \item \textbf{Manipulation (closed reduction):} a crooked nose may be repositioned under local or general anaesthesia, ideally within 7--14 days of injury
    \item \textbf{Septal haematoma:} urgent drainage by an ear, nose, and throat specialist to prevent necrosis and permanent damage to the septum
    \item \textbf{Surgery:} septorhinoplasty may be offered later for complex fractures, severe displacement, or persistent nasal obstruction
    \item \textbf{Activity advice:} avoid contact sports and further trauma to the nose for about six weeks, and avoid forceful nose-blowing in the first days
\end{itemize}

\section*{Complications}
\begin{itemize}
    \item Permanent change in the shape of the nose when a displaced fracture is not corrected
    \item A deviated septum causing chronic nasal obstruction
    \item Septal haematoma and, if untreated, septal infection, necrosis, or a perforation (hole) in the septum
    \item Recurrent or persistent nosebleeds
    \item Associated facial fractures involving the orbit, cheekbone, tear duct, or teeth
    \item Cerebrospinal fluid leak with a risk of meningitis after a skull-base fracture
    \item Residual cosmetic deformity that may require functional or cosmetic surgery later
\end{itemize}

\section*{Prognosis and Outcomes}
Most broken noses heal well without lasting problems. Swelling and bruising peak at two to three days and settle over one to two weeks, and the bones usually unite within three to six weeks. If the nose is clearly crooked, repositioning within the first two weeks gives the best cosmetic and functional result; delay makes correction more difficult and often requires surgery. A minority of people are left with a visible hump, a slight crook, or persistent nasal blockage and may consider a later septorhinoplasty. In children, continued facial growth can alter the final appearance of a healed nasal fracture, so paediatric injuries merit specialist review.

\section*{Prevention and Self-Care}
\begin{itemize}
    \item Wear appropriate protective headgear and face guards for contact and collision sports
    \item Wear helmets and seatbelts when cycling, motorcycling, or travelling in vehicles
    \item Avoid situations that increase the risk of fights and assaults, including heavy alcohol use
    \item After any nasal injury, apply ice early and keep the head elevated to limit swelling
    \item Avoid blowing the nose forcefully for the first few days after injury
    \item Avoid contact sports and rough play for around six weeks after a nasal fracture
    \item Attend follow-up so any deformity is assessed once the swelling has settled
\end{itemize}

\section*{Sources and Further Reading}
The following organisations provide reliable information on nasal fractures, first aid for facial injuries, and when to seek urgent care.

\begin{itemize}
    \item NHS. \textit{Nasal fracture (broken nose): symptoms, treatment, and first aid.} https://www.nhs.uk/conditions/broken-nose/
    \item Mayo Clinic. \textit{Broken nose: causes, treatment, and self-care.} https://www.mayoclinic.org/
    \item MedlinePlus. \textit{Nasal fracture and nose injury health information.} https://medlineplus.gov/
    \item MSD Manual. \textit{Fractures of the nose and facial bones.} https://www.msdmanuals.com/
    \item American Academy of Otolaryngology. \textit{Nasal fractures and facial trauma patient information.} https://www.enthealth.org/
    \item National Institute for Health and Care Excellence. \textit{Head injury and facial trauma guidance.} https://www.nice.org.uk/
\end{itemize}
